\subsection*{CU-18: Finalizar la partida}
\addcontentsline{toc}{subsection}{CU-18: Finalizar la partida}
\label{cu-18}

\begin{fichacu}{CU-18}{Finalizar la partida}
  \crudFila{Responsable}{Ángel Gabriel Aguilar Hernández}
  \crudFila{Actor principal}{Ninguno: caso incluido por \mbox{CU-13}, \mbox{CU-14}, \mbox{CU-15}, \mbox{CU-16}, \mbox{CU-17} y \mbox{CU-19}}
  \crudFila{Actores de sistema}{Servidor de partidas, Base de datos}
  \crudFila{Actores interesados}{Jugador}
  \crudFila{Prototipos}{2g, 2h, 15c, 15d, 17d}
  \crudFila{Prioridad}{Must (debe estar en la entrega). Categoría (b) Núcleo de jugabilidad. Caso incluido: condición de victoria, tiempo agotado y cierre de la partida con su elo y sus estadísticas.}
  \crudFila{Objetivo}{Cerrar la partida, determinar el resultado de cada participante y aplicar lo que la partida cambia en la cuenta: el elo y los contadores de la \textit{EstadisticaModo}, si la partida es clasificatoria.}
  \crudFila{Disparador}{Lo incluye el caso de uso que produce la condición de término: el \mbox{CU-13} cuando un peón llega a su fila de meta o el reloj se agota antes del movimiento, el \mbox{CU-14} cuando el reloj se agota antes de colocar el muro, el \mbox{CU-15} cuando se aceptan tablas, el \mbox{CU-16} cuando un jugador se rinde, el \mbox{CU-17} cuando un jugador abandona o su rival reclama la victoria, y el \mbox{CU-19} cuando el reloj de un jugador se agotó mientras estaba desconectado.}
  \textbf{Precondiciones} & \cuCelda{\begin{cureglas}
      \item \textbf{\mbox{PRE-01}} Existe una \textit{Partida} con estado \texttt{EN\_\allowbreak{}CURSO}.
      \item \textbf{\mbox{PRE-02}} Se ha producido una condición de término conforme a \mbox{RN-01}.
      \item \textbf{\mbox{PRE-03}} El Servidor de partidas tiene conexión disponible con la Base de datos.
    \end{cureglas}} \\ \hline
  \textbf{Flujo normal} & \cuCelda{\begin{cuenum}[start=1]
      \item El Servidor de partidas detecta la condición de término y determina la forma de término: \texttt{META}, \texttt{TIEMPO\_\allowbreak{}AGOTADO}, \texttt{RENDICION}, \texttt{ABANDONO} o \texttt{TABLAS} (\mbox{RN-01}). \textit{(Ver \mbox{FA-02}, \mbox{FA-03}, \mbox{FA-04})}
      \item El Servidor de partidas determina el resultado de cada \textit{Participacion}: ganada, perdida o tablas, y la posición en el modo de cuatro jugadores (\mbox{RN-02}).
      \item El Servidor de partidas comprueba, por el \texttt{tipo} de la \textit{Partida}, si la partida es clasificatoria, es decir, si nació del emparejamiento del \mbox{CU-09} (\mbox{RN-03}). \textit{(Ver \mbox{FA-01})}
      \item El Servidor de partidas inicia una transacción sobre la Base de datos.
      \item El Servidor de partidas actualiza la \textit{Partida}: estado \texttt{FINALIZADA}, \texttt{fecha\_\allowbreak{}fin} y forma de término. El ganador, si lo hay, no se escribe en la \textit{Partida}: es la \textit{Participacion} con resultado ganada. \textit{(Ver \mbox{EX-01})}
      \item El Servidor de partidas escribe en cada \textit{Participacion} su resultado, su posición y el tiempo de reloj que le quedaba.
      \item El Servidor de partidas calcula el elo nuevo de cada participante a partir del elo con el que entró, guardado en su \textit{Participacion}, y del resultado (\mbox{RN-04}). \textit{(Ver \mbox{EX-04})}
      \item El Servidor de partidas escribe en cada \textit{Participacion} el elo resultante, que junto con el de entrada da la diferencia que muestran la pantalla de fin y el historial (\mbox{RN-05}).
    \end{cuenum}} \\
   & \cuCelda{\begin{cuenum}[start=9]
      \item El Servidor de partidas actualiza la \textit{EstadisticaModo} de cada participante en el \textit{Modo} de la partida: incrementa \texttt{partidas\_\allowbreak{}jugadas}, incrementa \texttt{partidas\_\allowbreak{}ganadas} o \texttt{partidas\_\allowbreak{}perdidas}, escribe \texttt{puntos\_\allowbreak{}elo} y escribe \texttt{fecha\_\allowbreak{}ultima\_\allowbreak{}partida} con la \texttt{fecha\_\allowbreak{}fin} (\mbox{RN-06}). \textit{(Ver \mbox{EX-01})}
      \item Si la partida nació de una \textit{Sala}, el Servidor de partidas cambia su \texttt{estado} a \texttt{CERRADA} y elimina sus \textit{SalaParticipante}.
      \item El Servidor de partidas comprueba que ninguna \textit{Participacion} quede sin resultado: con él quedan terminadas, y sus jugadores vuelven a poder entrar a otra partida (\mbox{RN-07}).
      \item El Servidor de partidas confirma la transacción. \textit{(Ver \mbox{EX-02})}
      \item El Servidor de partidas notifica a cada cliente con \texttt{MATCH\_\allowbreak{}ENDED}, incluyendo el resultado, la forma de término y el cambio de elo. \textit{(Ver \mbox{FA-05}, \mbox{EX-03})}
      \item El SJM despliega la \texttt{GUI\_\allowbreak{}MatchEnd} con el resumen sobre el tablero, que sigue visible detrás.
      \item El SJM ofrece las opciones ``Buscar otra'', que continúa en el \mbox{CU-09} y no se ofrece a los invitados, ``Añadir al rival'', que tampoco se ofrece a los invitados porque no tienen amigos (\mbox{RN-07} del \mbox{CU-21}), y ``Volver al menú''.
      \item Fin del caso de uso.
    \end{cuenum}} \\ \hline
  \textbf{Flujos alternos} & \cuCelda{\textbf{\mbox{FA-01}: La partida no es clasificatoria}\par
    \begin{cuenum}[start=1]
      \item El Servidor de partidas omite los pasos 7, 8 y 9 del FN: no hay elo (\mbox{RN-03}).
      \item El Servidor de partidas no modifica la \textit{EstadisticaModo}: sus contadores describen solo partidas clasificatorias, para que cinco partidas privadas no basten para entrar en el ranking (\mbox{RN-12}).
      \item El flujo continúa en el paso 4 del FN.
    \end{cuenum}} \\
   & \cuCelda{\textbf{\mbox{FA-02}: La partida termina por abandono en el modo de cuatro jugadores}\par
    \begin{cuenum}[start=1]
      \item El Servidor de partidas marca la \textit{Participacion} del que abandonó con resultado perdido y forma de término \texttt{ABANDONO}, y retira su peón del tablero.
      \item Los muros que ese jugador había colocado permanecen en el tablero (\mbox{RN-08}).
      \item Si quedan al menos dos jugadores activos, la partida \textbf{no} termina: el flujo vuelve al \mbox{CU-13} y este caso solo se aplicó a esa participación. En caso contrario, el flujo continúa en el paso 1 del FN.
    \end{cuenum}} \\
   & \cuCelda{\textbf{\mbox{FA-03}: La partida termina en tablas}\par
    \begin{cuenum}[start=1]
      \item En el paso 2 del FN, el Servidor de partidas asigna el resultado de tablas a todas las \textit{Participacion}.
      \item En el paso 7, el Servidor de partidas calcula el elo conforme a \mbox{RN-04} tratando el resultado como intermedio, de modo que el de menor elo gane algo y el de mayor elo pierda algo.
      \item En el paso 9, la \textit{EstadisticaModo} incrementa \texttt{partidas\_\allowbreak{}jugadas} sin incrementar ganadas ni perdidas (\mbox{RN-06}).
      \item El flujo continúa en el paso 2 del FN.
    \end{cuenum}} \\
   & \cuCelda{\textbf{\mbox{FA-04}: La partida se gana por reloj}\par
    \begin{cuenum}[start=1]
      \item El Servidor de partidas registra la forma de término \texttt{TIEMPO\_\allowbreak{}AGOTADO}, que distingue esta victoria de una por llegada a meta (\mbox{RN-09}).
      \item El flujo continúa en el paso 2 del FN.
    \end{cuenum}} \\
   & \cuCelda{\textbf{\mbox{FA-05}: Un participante tiene la sesión cerrada al terminar la partida}\par
    \begin{cuenum}[start=1]
      \item El Servidor de partidas aplica igualmente todos los pasos del FN sobre ese participante: la progresión no depende de que esté conectado.
      \item El Servidor de partidas no le envía \texttt{MATCH\_\allowbreak{}ENDED}.
      \item El jugador verá el resultado en su historial la próxima vez que entre.
      \item El flujo continúa en el paso 14 del FN para los demás.
    \end{cuenum}} \\ \hline
  \textbf{Excepciones} & \cuCelda{\textbf{\mbox{EX-01}: Fallo al escribir en la base de datos}\par
    \begin{cuenum}[start=1]
      \item El Servidor de partidas revierte la transacción completa, de modo que no quede una partida cerrada con estadísticas a medio actualizar (\mbox{RN-10}).
      \item El Servidor de partidas registra la excepción en la bitácora del servidor con un identificador de incidencia.
      \item El Servidor de partidas deja la \textit{Partida} en \texttt{EN\_\allowbreak{}CURSO} y reintenta el cierre conforme a su política de reintentos.
      \item Si los reintentos se agotan, el Servidor de partidas marca la \textit{Partida} como \texttt{PENDIENTE\_\allowbreak{}DE\_\allowbreak{}CIERRE} y notifica a los clientes que el resultado se aplicará más tarde (\mbox{RN-11}).
      \item Fin del caso de uso.
    \end{cuenum}} \\
   & \cuCelda{\textbf{\mbox{EX-02}: Fallo al confirmar la transacción}\par
    \begin{cuenum}[start=1]
      \item El Servidor de partidas trata la transacción como fallida y no notifica el fin a nadie.
      \item El Servidor de partidas registra la excepción señalando que el estado del cierre es indeterminado.
      \item El flujo continúa en el paso 3 del \mbox{EX-01}.
    \end{cuenum}} \\
   & \cuCelda{\textbf{\mbox{EX-03}: Un cliente no recibe la notificación de fin}\par
    \begin{cuenum}[start=1]
      \item El Servidor de partidas no reintenta indefinidamente: el resultado ya está escrito y es consultable.
      \item Al reconectar, el SJM detecta que su \textit{Participacion} terminó y despliega la \texttt{GUI\_\allowbreak{}MatchEnd} con el resultado leído del servidor.
      \item Fin del caso de uso.
    \end{cuenum}} \\
   & \cuCelda{\textbf{\mbox{EX-04}: El cálculo de elo no puede completarse}\par
    \begin{cuenum}[start=1]
      \item El Servidor de partidas detecta que falta el elo inicial en alguna \textit{Participacion}, lo que solo puede ocurrir por una partida creada de forma incorrecta.
      \item El Servidor de partidas cierra la partida sin aplicar elo, y registra la incidencia en la bitácora del servidor para su revisión.
      \item El flujo continúa en el paso 10 del FN.
    \end{cuenum}} \\ \hline
  \textbf{Postcondiciones} & \cuCelda{\begin{cureglas}
      \item \textbf{\mbox{POST-01}} La \textit{Partida} tiene estado \texttt{FINALIZADA}, su \texttt{fecha\_\allowbreak{}fin}, su forma de término; su ganador, si lo hay, es la \textit{Participacion} con resultado ganada.
      \item \textbf{\mbox{POST-02}} Cada \textit{Participacion} tiene su resultado y su posición y, si la partida era clasificatoria, su elo final y la diferencia con el inicial.
      \item \textbf{\mbox{POST-03}} Ninguna \textit{Participacion} de la partida sigue sin terminar, de modo que todos sus jugadores pueden entrar a otra partida (\mbox{D-01}).
      \item \textbf{\mbox{POST-04}} Si la partida era clasificatoria, la \textit{EstadisticaModo} de cada participante en ese \textit{Modo} refleja la partida en todos sus contadores.
      \item \textbf{\mbox{POST-05}} Si la partida no era clasificatoria, ningún elo cambió y ninguna \textit{EstadisticaModo} se modificó.
    \end{cureglas}} \\
   & \cuCelda{\begin{cureglas}
      \item \textbf{\mbox{POST-06}} Si el caso termina por una excepción, la partida no quedó cerrada a medias: o se aplicó todo, o sigue \texttt{EN\_\allowbreak{}CURSO} o \texttt{PENDIENTE\_\allowbreak{}DE\_\allowbreak{}CIERRE}.
      \item \textbf{\mbox{POST-07}} Las \textit{Jugada} de la partida siguen íntegras: el cierre no las modifica.
      \item \textbf{\mbox{POST-08}} Si la partida quedó \texttt{FINALIZADA} y nació de una \textit{Sala}, esa sala está \texttt{CERRADA} y sin miembros.
    \end{cureglas}} \\ \hline
  \textbf{Reglas de negocio} & \cuCelda{\begin{cureglas}
      \item \textbf{\mbox{RN-01}} Una partida termina de cinco formas: llegada a meta, reloj agotado, rendición, tablas aceptadas y abandono. La forma de término se guarda, no solo el resultado, porque no todas valen lo mismo.
      \item \textbf{\mbox{RN-02}} En el modo de cuatro jugadores, quien llega a meta queda clasificado con su posición y la partida sigue para los demás.
      \item \textbf{\mbox{RN-03}} Solo las partidas clasificatorias, las que crea el emparejamiento del \mbox{CU-09} con \texttt{tipo} \texttt{CLASIFICATORIA}, cambian el elo y suman a la \textit{EstadisticaModo}, que alimenta el ranking. Las de sala privada, con \texttt{tipo} \texttt{PRIVADA}, son amistosas: no cambian el elo ni la \textit{EstadisticaModo}.
      \item \textbf{\mbox{RN-04}} El elo se calcula a partir del elo de entrada de cada participante, no del que tenga en ese momento en su cuenta. Dos partidas simultáneas no pueden ocurrir (\mbox{D-01}), pero el elo de entrada es de todos modos lo único que hace reproducible el cálculo.
      \item \textbf{\mbox{RN-05}} El cambio de elo se guarda en la \textit{Participacion}. El historial del \mbox{CU-26} lo muestra partida por partida y no podría recomponerlo restando estados sucesivos de la cuenta.
      \item \textbf{\mbox{RN-06}} Los contadores de \textit{EstadisticaModo} son por modo, no globales. El elo del menú y del ranking general es el del modo clásico 9×9.
    \end{cureglas}} \\
   & \cuCelda{\begin{cureglas}
      \item \textbf{\mbox{RN-07}} Cerrar la partida cierra todas sus participaciones. Mientras una siga sin terminar, su jugador no puede entrar a otra partida.
      \item \textbf{\mbox{RN-08}} Los muros de un jugador que abandonó permanecen en el tablero. El abandono retira al jugador, no lo que ya hizo.
      \item \textbf{\mbox{RN-09}} Una victoria por reloj se distingue de una victoria por llegada a meta.
      \item \textbf{\mbox{RN-10}} El cierre de la partida, la escritura de resultados, la actualización de estadísticas y, si la partida nació de una sala, el cierre de la sala ocurren en una sola transacción.
      \item \textbf{\mbox{RN-11}} Una partida que no puede cerrarse queda marcada como pendiente de cierre y se reintenta. No se descarta: los jugadores quedarían bloqueados para siempre por el \mbox{RN-07}.
      \item \textbf{\mbox{RN-12}} Hacen falta cinco partidas en un modo para aparecer en su clasificación.
    \end{cureglas}} \\ \hline
  \textbf{Observación} & \cuCelda{\textit{Sobre por qué este caso es un caso incluido.}\par
    Ningún actor lo inicia: lo incluyen los seis casos que pueden terminar una partida, el \mbox{CU-13} por llegada a meta o reloj agotado, el \mbox{CU-14} por reloj agotado, el \mbox{CU-15} por tablas, el \mbox{CU-16} por rendición, el \mbox{CU-17} por abandono y el \mbox{CU-19} por reloj agotado durante una desconexión. Separarlo concentra en un solo lugar las reglas de victoria y la actualización del elo y de las estadísticas que los seis comparten; si cada uno las repitiera, bastaría corregir una regla en un caso y olvidarla en otro para que la misma partida puntuara distinto según cómo terminó.} \\ \hline
  \textbf{Observación} & \cuCelda{\textit{Sobre por qué este caso es el único que escribe el elo.}\par
    Es el caso de uso más caro del sistema y el que más tablas toca de una vez, y por eso conviene decir en voz alta por qué no se reparte. El resultado, el elo y los contadores dependen todos del mismo cierre y deben moverse juntos o no moverse: un jugador con el elo actualizado y los contadores sin tocar es un error que nadie sabría reparar después. Repartirlo entre varios casos —uno que cierre, otro que puntúe, otro que cuente— obligaría a coordinarlos con una transacción que los abarcara a todos, que es exactamente lo que este caso ya es.} \\ \hline
  \textbf{Datos que toca} & \cuCelda{\begin{cudatos}
      \item \textbf{\textit{Partida}:} lee \texttt{tipo}; escribe estado, \texttt{fecha\_\allowbreak{}fin}, forma de término \textit{(pasos 3, 5)}
      \item \textbf{\textit{Participacion}:} escribe resultado, posición, elo final, reloj restante \textit{(pasos 6, 8, 11)}
      \item \textbf{\textit{EstadisticaModo}:} escribe \texttt{partidas\_\allowbreak{}jugadas}, \texttt{partidas\_\allowbreak{}ganadas}, \texttt{partidas\_\allowbreak{}perdidas}, \texttt{puntos\_\allowbreak{}elo} y \texttt{fecha\_\allowbreak{}ultima\_\allowbreak{}partida} \textit{(paso 9)}
      \item \textbf{\textit{Sala}, \textit{SalaParticipante}:} cierra la sala de origen, si la hay, y elimina sus miembros \textit{(paso 10)}
      \item \textbf{\textit{Modo}:} lee, para los parámetros de puntuación \textit{(paso 7)}
    \end{cudatos}} \\ \hline
\end{fichacu}
\clearpage
